% part: intuitionistic-logic
% chapter: introduction
% section: constructive-reasoning

\documentclass[../../../include/open-logic-chapter]{subfiles}

\begin{document}

\olfileid{int}{int}{cr}

\section{تعمیری دلیل}

امکانی عاملاں یا دوجے درجے دے مسوّراں نال کلاسیکی منطق ودھ جاندی
اے۔ وجدانیت پسند منطق کلاسیکی منطق اُتے قید لاؤندی اے، تے ایس لحاظ
نال «غیر کلاسیکی» اے۔ کلاسیکی منطق کئی طرحاں نال \emph{غیر تعمیری} اے۔
وجدانیت پسند منطق دا مقصد اوس تعمیری دلیل نوں رسمی شکل دینا اے جیہڑی
اک قسم دی تعمیری ریاضی وچ ورتی جاندی اے۔ ہِیٹھلیاں مثالاں ایس پچھے
کم کردیاں کجھ وجھاں سمجھا سکدیاں نیں۔

منّ لو کوئی آکھے کہ اوہنے اک قدرتی عدد~$n$ لبھیا اے جیہدی خاصیت
ایہ اے کہ جے $n$ جفت اے، تاں ریمن دا مفروضہ سچ اے، تے جے $n$ طاق
اے، تاں ریمن دا مفروضہ جھوٹ اے۔ بڑی چنگی خبر!{} ریمن دا مفروضہ سچ
اے یا نہیں، ایہ ریاضی دے وڈے اَن حل سوالاں وچوں اک اے۔ لگدا اے
اوہنے ایس مسئلے نوں اک حساب تک لے آندا اے، یعنی صرف ایہ ویکھنا
رہ گیا اے کہ خاص عدد جفت اے یا نہیں۔

~$n$ دی ایہ کراماتی قدر کیہڑی اے؟ اوہ ایویں دسدا اے کہ $n$ اوہ
قدرتی عدد اے جیہڑا $2$ دے برابر اے جے ریمن دا مفروضہ سچ اے، تے
ہور صورت وچ $3$ دے برابر اے۔

تسی غصے وچ اپنے پیسے واپس منگدے او۔ کلاسیکی نظر توں، اُتلا بیان
سچی مچی $n$ دی اکّو قدر متعین کردا اے۔ پر تہانوں $n$ دی اوہ قدر
چاہیدی اے جیہڑی \emph{صاف صاف} دتی ہووے۔

اک ہور مثال ویکھو جیہڑی شاید گھٹ بناوٹی اے۔ سانوں پتا اے کہ
غیر ناطق عدد نوں ناطق قوّت تائیں چُک کے ناطق نتیجہ آ سکدا اے۔
مثلاً $\sqrt{2}^2 = 2$۔ پر ایہ گل اینی صاف نہیں کہ غیر ناطق عدد
نوں \emph{غیر ناطق} قوّت تائیں چُک کے ناطق نتیجہ آ سکدا اے یا نہیں۔
ہِیٹھلا قضیہ دسدا اے کہ ایہ ہو سکدا اے:

\begin{thm}
ایسے غیر ناطق عدد $a$ تے $b$ موجود نیں جیہناں لئی $a^b$ ناطق اے۔
\end{thm}

\begin{proof}
$\sqrt{2}^{\sqrt{2}}$ نوں ویکھو۔ جے ایہ ناطق اے، تاں کم ہو گیا۔
اسیں $a = b = \sqrt{2}$ رکھ سکدے آں۔ ہور صورت وچ ایہ غیر ناطق
اے۔ فیر ساڈے کول
\[
(\sqrt{2}^{\sqrt{2}})^{\sqrt{2}} = \sqrt{2}^{\sqrt{2} \cdot
  \sqrt{2}} = \sqrt{2}^2 = 2,
\]
اے، جیہڑا ناطق اے۔ سو ایس صورت وچ $a$ نوں $\sqrt{2}^{\sqrt{2}}$
تے $b$ نوں~$\sqrt 2$ رکھو۔
\end{proof}

کی ایہ معتبر ثبوت اے؟ زیادہ تر ریاضی دان اینوں معتبر سمجھدے
نیں۔ پر ایتھے وی کجھ کھٹکدا اے۔ اسیں ثابت کیتا اے کہ خاص خاصیت
والی حقیقی عدداں دی جوڑی موجود اے، پر ایہ نہیں دس سکدے کہ
اوہ \emph{کیہڑی} جوڑی اے۔ ایہو نتیجہ ایویں وی ثابت ہو سکدا اے
کہ جوڑی $a$، $b$ ثبوت وچ \emph{دتی ہووے}۔ $a = \sqrt{3}$ تے
$b = \log_3 4$ رکھو۔ فیر
\[
a^b = \sqrt{3}^{\log_3 4} = 3^{1/2 \cdot \log_3 4} = (3^{\log_3
  4})^{1/2} = 4^{1/2}= 2,
\]
کیوں جے $3^{\log_3 x} = x$۔

وجدانیت پسند منطق اوس طرح دی دلیل نوں رسمی شکل دیندی اے جیہدے
وچ پہلے ثبوت ورگے قدم دی اجازت نہیں۔ ایہ ثابت کرن لئی کہ کوئی $x$
~$!A(x)$ نوں پورا کردا اے، اک خاص~$x$ دینا پیندا اے، تے ایہدا
ثبوت وی کہ اوہ $!A$ نوں پورا کردا اے، جیویں دوجے ثبوت وچ اے۔
$!A$ یا $!B$ دے لاگو ہون دا ثبوت دین لئی دونہاں وچوں اک دا ثبوت
دے سکنا چاہیدا اے۔

رسمی لحاظ نال، کلاسیکی منطق دے کسے !!a{derivation} نظام اُتے
اک خاص طرح دی قید لاؤن نال وجدانیت پسند منطق ملدی اے۔ ریاضی دی
نظر توں ایہ رسمی اخذ نظام نیں۔ پر جیویں پہلاں آکھیا اے، ایہناں
دا مقصد ریاضیاتی دلیل دی اک قسم نوں رسمی شکل دینا اے۔ کوئی
اینوں ریاضی دے کسے فلسفیانہ خیال تحت جائز دلیل سمجھ سکدا اے،
مثلاً براوَر دی وجدانیت۔ کوئی اینوں ودھ «ٹھوس» تے تسلی بخش
ریاضیاتی دلیل سمجھ سکدا اے، جیویں بِشپ دی تعمیریت وچ۔ ایس گل
اُتے بحث ہو سکدی اے کہ رسمی بیان پچھے کم کردی غیر رسمی وجھ نوں
پھڑدا اے یا نہیں۔ ساڈی فلسفیانہ سوچ بھاویں جیہڑی وی ہووے، اسیں
وجدانیت پسند منطق نوں رسمی منطق دے طور اُتے پڑھ سکدے آں۔ کئی
ریاضیاتی منطق دان ایہ پڑھائی وکھ وکھ وجھاں کرکے دلچسپ سمجھدے نیں۔

\end{document}
